
\section*{Alban au temps de la réforme (1562-1598)}

\lettrine{L}{'}agitation religieuse commence d'assez bonne heure dans
notre montagne. Ce n'est pas que les « Albans » soient plus combatifs
que d'autres, mais, selon toute apparence, les querelles qui
s'engagent sont assez souvent le fait d'hommes sans foi ni loi.
La dévotion n'est qu'un prétexte pour terroriser un pays, piller et
massacrer de braves gens au nom du Christ. Aussi les habitants
finissent par donner à ces bandes de pillards, qu'elles s'étiquettent
catholiques ou protestantes, le nom général de « voleurs », de
brigands. Ne faisons pas trop les censeurs de ces temps révolus,
on sait hélas, qu'à l'heure actuelle et au nom d'idéologies
différentes, on se massacre à qui mieux mieux et avec des moyens
tels qui font que les guerres d'antan nous paraissent jeux d'enfants
à côté des tueries de ce XX\textsuperscript{ème} siècle agonisant.

\lettrine{C}{'}est quelques unes de ces bagarres dites de
« religion » qui vont faire les frais de cette petite chronique
locale. On raconte qu'en 1557, c'était au château des Farguettes
près Crespinet, qu'un Albanais, Jean Puech, se livra à une plaisanterie
de mauvais goût. Une famille Puech est signalée à Ginestous dans
le cadastre de 1598. Et il y eut des Puech au Noyer jusqu'au début
du XX\textsuperscript{ème} siècle. Une croix érigée par eux existe
encore sur le petit chemin de traverse qui unit la voie d'Alban au
Noyer et celle du Noyer à Ginestous. Notre compatriote avec un
villefranchois : Jean Albert qui se titrera un jour de seigneur
baron d'Alban, s'en donne à cœur joie pour parodier une cérémonie
d'église. Nos deux gaillards ont donc, pour la circonstance, revêtu
des oripeaux divers et fait jouer à une ânesse le rôle de pénitente
à l'article de la mort. Entre joyeux lurons, quelque peu éméchés,
cette mascarade ne mérite pas les
